\section{Requisitos do Usuário}

Os requisitos de usuário descrevem as necessidades que motivam o
desenvolvimento do sistema e os serviços que o usuário espera obter.
Nesta etapa, a descrição é feita de forma compreensível ao usuário,
sem antecipar detalhes de implementação.

\subsection{Visão do Usuário}

O usuário precisa de uma ferramenta que o auxilie a analisar ativos
financeiros e a montar diferentes carteiras de investimentos,
considerando a relação entre o risco que está disposto a assumir e o
retorno esperado.

O sistema deverá permitir a simulação de diferentes distribuições
de capital entre os ativos selecionados e apresentar os resultados
de forma clara, possibilitando a identificação das combinações
eficientes segundo o modelo de Markowitz.

\subsection{O que o Usuário Espera Fazer}

O usuário deverá ser capaz de:

\begin{itemize}
\item selecionar os ativos que participarão da análise;

\item consultar a lista de empresas listadas na bolsa, com seus
respectivos códigos (tickers), para seleção dos ativos;\textbf{}

\item obter automaticamente, por meio de uma API, os dados
históricos necessários para a análise;

\item fornecer ou carregar os dados históricos necessários para
a análise;

\item simular diferentes combinações de ativos e pesos;

\item consultar o retorno esperado de cada carteira analisada;

\item consultar o risco associado a cada carteira;

\item comparar diferentes composições de carteira;

\item identificar as carteiras pertencentes à fronteira
eficiente;

\item consultar a participação de cada ativo nas carteiras
analisadas;

\item visualizar os resultados por meio de tabelas e gráficos.


\end{itemize}

\subsection{Resultado Esperado}

Ao utilizar o sistema, o usuário deverá obter informações que
permitam comparar diferentes alternativas de composição de carteira
e avaliar a relação entre risco e retorno.

Dessa forma, o sistema deverá auxiliar o usuário a identificar
carteiras eficientes segundo os critérios estabelecidos pelo modelo
de Markowitz.
